% !TeX root = ../../arithmetic-and-computation.tex
\section{Historical exclusion batches and their exact scope}
The following records retain their original target populations and counts.
Only the complete two-region theorem of Paper I supplies the full-curve
conclusion for the whole packet. In every statement below that concludes
curve exclusion, the packet execution premise of Paper I is inherited;
no local summary removes that premise.

\paragraph{Historical first-target computations and their present use.}
The following subpacket tables retain the original-product $P_0$ counts,
matrix lists and numerical mask digests. Their old checkpoint counts
such as $107$, $105$ and $75$ describe that historical progression, not
the current undecided population. The full-curve exclusion statements
use Theorem~\externalref{packet}{thm:fullpacketgeometry} in addition to these computations.
For the nonprincipal proper-$R_{83}$ module, the old and current
arithmetic censuses each describe the entire same $32$-point population.
The new reflected search excludes all of it. No permutation identifying
an individual old numerical mask with a current mask is assumed or
needed; no historical digest is relabelled as a reflected digest.

\begin{proposition}[Two proper-order $d=7$ exclusions]\label{prop:d7mask14}
In the canonical point manifest of Theorem~\ref{thm:d7ambient}, the
two nontrivial-$R_{83}$ points represented by the free root pairs
$\{14,62\}$ and $\{30,46\}$ do not lie on the normalized McMullen curve.
In the historical v45 first-target bookkeeping, this pair reduced $107$ to $105$:
$11$ maximal-multiplier and $94$ proper-order points.
\end{proposition}
\begin{proof}
Take the index-$2$ action on the nontrivial conductor-$83$ ideal lattice used
in Theorem~\ref{thm:d7ambient}, with determinant-sector mask $14$.  The
absolute norm of its raw lower-left matrix entry is
\[
2097130411225081.
\]
Thus diagonal unit balancing alone preserves a prohibitively small product of
the six coordinate heights.  The integral determinant-one
inversion
\[
 S=\begin{pmatrix}0&-1\\1&0\end{pmatrix}
\]
replaces it by the negative upper-right entry, of absolute norm $83^2=6889$.
Put $u=4-t^2$ and then apply
$H=\operatorname{diag}(u^{-2},u^2)\in\mathrm{SL}_2(\OL)$.
Both transformations have integral inverses and hence preserve the orbit
membership question.  For the resulting point the exact squared-height norm
is $1/28217344$.

The compact polygon bounds used in
Proposition~\ref{prop:residualexclusions} are rederived for this source point,
using its exact $\operatorname{Re}w$, $\operatorname{Im}(w)^2$, and $|w|^2$.
Outward rational bounds and the positive trace form give exactly $23923$
possible $a$-entries, $6605$ possible $c$-entries, and $2093$ coupled bottom
rows $(c,d)$.  For $c\ne0$, the determinant equation forces
$b=(ad-1)/c$; exact divisibility and numerator bounds leave $9036$ matrices.
The separate $c=0$ branch is empty.  Exact cross-multiplied circle signs reject
every matrix from at least one coordinate quadrilateral, while retaining
boundary equality.  Thus there is no polygon survivor.

For the second point, bit $16$ changes the orientation unit by $-1$.  If
$J=\operatorname{diag}(1,-1)$, exact matrix arithmetic gives
$M_{30}=JM_{14}J^{-1}$ for the reduced matrices, hence their fixed points have
centers $x$ and $-x$ and common squared height and modulus.  Here
$\det J=-1$ is totally negative, so $J$ lies in neither
$\GL_2^+(\OL)$ nor $\SL_2(\OL)$ and does not identify the quotient points.
Nevertheless
\[
 \iota\begin{pmatrix}a&b\\c&d\end{pmatrix}
 =\begin{pmatrix}a&-b\\-c&d\end{pmatrix}
\]
is an involution of $\SL_2(\OL)$ that bijects their complete bounded matrix
sets, including the $c=0$ sector: for $j(z)=-\overline z$ one has
$\iota(\gamma)j(z)=j(\gamma z)$.  We do not assume that the quadrilateral is
reflection-invariant: every reflected matrix is tested separately at the
mask-$30$ point.

The primary rational/SymPy and separate standard-library/Sturm implementations
reproduce both empty survivor sets and the respective canonical exclusion
digests
\[
\begin{gathered}
\digest{19f1bba6f56200bac45a52810a56669c}{51824f22a9b3fd410db6a1fc2d69a5a1},\\
\digest{2334a40fed0d008d44c2e892871bec20}{fc04b226a99df899a8927df81595faf6}.
\end{gathered}
\]
The root involution is xor $48$, so $\{14,62\}$ and $\{30,46\}$ are disjoint
points.  Removing them from the $107$ points left by
Theorem~\ref{thm:d7ambient} gives $11+94=105$.
\end{proof}


\begin{proposition}[Two further proper-order $d=7$ exclusions]\label{prop:d7mask15_31}
In the canonical point manifest of Theorem~\ref{thm:d7ambient}, the
two nontrivial-$R_{83}$ points represented by the free root pairs
$\{15,63\}$ and $\{31,47\}$ do not lie on the normalized McMullen curve.
In the historical v47 first-target bookkeeping, this pair reduced $105$ to $103$:
$11$ maximal-multiplier and $92$ proper-order points.
\end{proposition}
\begin{proof}
Use the index-$2$ action on the nontrivial conductor-$83$ ideal lattice.  For
each of masks $15$ and $31$, first apply the integral determinant-one inversion
from Proposition~\ref{prop:d7mask14}, then conjugate by
$\operatorname{diag}(q,q^{-1})$, where $q$ has saturated-unit exponent vector
$(1,1,1,0,0)$ in $(\xi_5,\xi_{11},\xi_{13},\xi_{17},4-t^2)$.  Both maps have
integral inverses.  The resulting lower-left entry has absolute norm $83^2$,
and the exact squared-height norm is $1/28217344$.

The point-specific compactness bounds are rederived with rational Sturm
intervals and rounded outward.  Complete positive-trace-form enumeration gives
$24103$ possible $a$-entries, $6611$ possible $c$-entries, and $2045$ coupled
bottom rows.  Exact determinant solving and coefficient divisibility leave
$8278$ matrices with $c\ne0$; the separately enumerated $c=0$ branch is empty.

For $J=\operatorname{diag}(1,-1)$, the two reduced action matrices satisfy
$M_{31}=JM_{15}J^{-1}$.  The determinant of $J$ is totally negative, so
$J\in\GL_2(\OL)\setminus\GL_2^+(\OL)$ and $J\notin\SL_2(\OL)$: this does not
identify the two quotient points.  As in Proposition~\ref{prop:d7mask14}, the
involution
\[
 \begin{pmatrix}a&b\\c&d\end{pmatrix}
 \longmapsto
 \begin{pmatrix}a&-b\\-c&d\end{pmatrix}
\]
bijects their complete bounded determinant-one sets, in both directions and
including the $c=0$ sector.  No quadrilateral-reflection symmetry is assumed;
all four exact side predicates at all six real embeddings are evaluated
separately for both centers.  Neither point has a survivor.

The primary rational/SymPy and separate standard-library/Sturm implementations
reproduce the two canonical exclusion digests
\[
\begin{gathered}
\digest{66684a7a1aa493a7e37b6b4fc6e97b36}{63c70e79700724eaca94e170a4d946b4},\\
\digest{2fa340baeab5f877f4ea61d9c2e9cc85}{0ee54221d257396b1320dc2c048ef023}.
\end{gathered}
\]
The v45 root involution is xor $48$, so $\{15,63\}$ and $\{31,47\}$ are
disjoint and are also disjoint from the two fibers in
Proposition~\ref{prop:d7mask14}.  Removing these two further points from the
$105$ left there gives $11+92=103$.
\end{proof}


\begin{proposition}[A third pair of proper-order $d=7$ exclusions]\label{prop:d7mask6_22}
In the canonical point manifest of Theorem~\ref{thm:d7ambient}, the
two nontrivial-$R_{83}$ points represented by the free root pairs
$\{6,54\}$ and $\{22,38\}$ do not lie on the normalized McMullen curve.
At the historical v49 checkpoint, exactly $101$ ambient $d=7$ points
remained unclassified for curve membership: $11$ maximal-multiplier and $90$ proper-order points.
\end{proposition}
\begin{proof}
Apply the integral determinant-one inversion of
Proposition~\ref{prop:d7mask14} to the index-$2$ nontrivial conductor-$83$
ideal action, then conjugate by $\operatorname{diag}(q,q^{-1})$ with
$q=\xi_{11}\xi_{13}$.  This is the saturated-unit exponent vector
$(0,1,1,0,0)$. Both maps and their inverses are integral. The reduced
lower-left entry has absolute norm $83^2$, and the exact squared-height norm is
$1/28217344$.

Point-specific outward rational bounds, certified from Sturm intervals, give
$23937$ possible $a$-entries, $6605$ possible $c$-entries, and $2019$ coupled
bottom rows. Exact determinant solving and coefficient divisibility leave
$8514$ matrices with $c\ne0$; exhaustive enumeration of the separate $c=0$
branch gives none.

The reduced action matrices satisfy $M_{22}=JM_6J^{-1}$ for
$J=\operatorname{diag}(1,-1)$. Since $\det J=-1$ is totally negative,
$J\in\GL_2(\OL)\setminus\GL_2^+(\OL)$ and $J\notin\SL_2(\OL)$; it does not
identify the quotient points. The involution
\[
 \begin{pmatrix}a&b\\c&d\end{pmatrix}
 \longmapsto
 \begin{pmatrix}a&-b\\-c&d\end{pmatrix}
\]
bijects their complete bounded determinant-one sets in both directions,
including $c=0$, because the exact numerator and denominator norms agree at
the reflected centers. No quadrilateral-reflection symmetry is assumed: all
four exact side predicates at all six real embeddings are evaluated separately
at both points. Neither test has a survivor.

The primary rational/SymPy and separate standard-library/Sturm implementations
reproduce the canonical exclusion digests
\[
\begin{gathered}
\digest{756c33defcade60c47adab07c30af7b5}{beb7b110e2051e8dfe025055690cb523},\\
\digest{5f5c2adaa1366bceab1d2b43a49fffc4}{8c8377320e9d4d522853e1586950e04a}.
\end{gathered}
\]
The xor-$48$ fibers $\{6,54\}$ and $\{22,38\}$ are disjoint from one another
and from the four earlier excluded proper-order fibers. Removing these two
further points from the $103$ left by Proposition~\ref{prop:d7mask15_31}
gives $11+90=101$.
\end{proof}


\begin{proposition}[A fourth pair of proper-order $d=7$ exclusions]\label{prop:d7mask1_17}
In the canonical point manifest of Theorem~\ref{thm:d7ambient}, the
two nontrivial-$R_{83}$ points represented by the free root pairs
$\{1,49\}$ and $\{17,33\}$ do not lie on the normalized McMullen curve.
At the historical v50 checkpoint, exactly $99$ ambient $d=7$ points
remained unclassified for curve membership: $11$ maximal-multiplier and $88$ proper-order points.
\end{proposition}
\begin{proof}
Apply the integral determinant-one inversion of
Proposition~\ref{prop:d7mask14} to the index-$2$ nontrivial conductor-$83$
ideal action, then conjugate by $\operatorname{diag}(q,q^{-1})$ with
$q=\xi_5\xi_{17}^{-1}(4-t^2)^{-1}$. This is the saturated-unit exponent
vector $(1,0,0,-1,-1)$. Both maps and their inverses are integral. The reduced
lower-left entry has absolute norm $83^2$, and the exact squared-height norm is
$1/28217344$.

Point-specific outward rational bounds, certified from Sturm intervals, give
$23959$ possible $a$-entries, $6601$ possible $c$-entries, and $2017$ coupled
bottom rows. Exact determinant solving and coefficient divisibility leave
$8748$ matrices with $c\ne0$; exhaustive enumeration of the separate $c=0$
branch gives none.

The reduced action matrices satisfy $M_{17}=JM_1J^{-1}$ for
$J=\operatorname{diag}(1,-1)$. Since $\det J=-1$ is totally negative,
$J\in\GL_2(\OL)\setminus\GL_2^+(\OL)$ and $J\notin\SL_2(\OL)$; it does not
identify the quotient points. The involution
\[
 \begin{pmatrix}a&b\\c&d\end{pmatrix}
 \longmapsto
 \begin{pmatrix}a&-b\\-c&d\end{pmatrix}
\]
bijects their complete bounded determinant-one sets in both directions,
including $c=0$, because the exact numerator and denominator norms agree at
the reflected centers. No quadrilateral-reflection symmetry is assumed: all
four exact side predicates at all six real embeddings are evaluated separately
at both points. Neither test has a survivor.

The primary rational/SymPy and separate standard-library/Sturm implementations
reproduce the canonical exclusion digests
\[
\begin{gathered}
\digest{3d6a732cc36a03fe90933cdfcffd1b9d}{977e7ce22cb1d7ae1f3926e6dff2e1a8},\\
\digest{82b2d957ba62af939e5bb39b0f6bbec0}{6b230094aeee60faa3ae54caf3e2e17b}.
\end{gathered}
\]
The xor-$48$ fibers $\{1,49\}$ and $\{17,33\}$ are disjoint from one another
and from the six earlier excluded proper-order fibers. Removing these two
further points from the $101$ left by Proposition~\ref{prop:d7mask6_22}
gives $11+88=99$.
\end{proof}


\begin{proposition}[A twenty-point proper-order $d=7$ exclusion batch]\label{prop:d7batch10}
In the canonical point manifest of Theorem~\ref{thm:d7ambient}, the twenty
nontrivial-$R_{83}$ points with canonical masks
\[
 \{0,16,2,18,3,19,4,20,5,21,8,24,9,25,10,26,11,27,12,28\}
\]
do not lie on the normalized McMullen curve.  Here the point with mask $m$
is represented by the free root fiber $\{m,m\mathbin{\mathrm{xor}}48\}$.
At the historical v51 checkpoint, exactly $79$ ambient $d=7$ points
remained unclassified for curve membership: $11$ maximal-multiplier and $68$ proper-order points.
\end{proposition}
\begin{proof}
For each base mask
\[
 m\in\{0,2,3,4,5,8,9,10,11,12\},
\]
reconstruct the index-$2$ action on the nontrivial conductor-$83$ ideal
lattice. Apply the integral determinant-one inversion of
Proposition~\ref{prop:d7mask14}, followed by the point-specific saturated-unit
conjugation $\operatorname{diag}(q_m,q_m^{-1})$. Exact inversion in
$\OL$ verifies that $q_m$ and $q_m^{-1}$ are integral, so these reductions and
their inverses lie in $\SL_2(\OL)$. The reduced lower-left entries have
absolute norm $83^2$, while the exact squared-height and squared-modulus norms
are respectively
\[
 \frac1{28217344},\qquad
 \frac{2097130411225081}{6889}.
\]

For each base point, rational Sturm intervals derive point-specific outward
compact bounds. The exact positive trace form and its rational
$LDL^{\mathsf T}$ decomposition enumerate every algebraic integer in the
resulting six-place boxes. The complete counts are
\[
\begin{array}{c|rrrrr}
m&\#a&\#c&\#(c,d)&c\ne0&c=0\\ \hline
0 &24035&6633&2023&8136&0\\
2 &24137&6607&1967&8334&0\\
3 &23993&6597&1937&7794&0\\
4 &24059&6599&2015&8702&10\\
5 &23959&6619&2083&8878&0\\
8 &23967&6613&2021&8750&6\\
9 &24001&6607&2085&9146&0\\
10&24039&6611&2017&9436&6\\
11&23973&6617&2013&8360&0\\
12&24013&6635&2015&8508&4
\end{array}
\]
For $c\ne0$, the determinant equation forces $b=(ad-1)/c$, whose
integrality and numerator bound are decided exactly. For $c=0$, it forces
$a=d^{-1}$, after which a second exact trace-form enumeration exhausts every
bounded $b$. Thus the last two columns include the entire determinant-one
search. At mask $3$, the vector $(-7,0,0,0,0,0)$ meets the embedding-$1$
$c$-entry bound exactly; the closed-box comparison accepts that local
equality, while another embedding rejects the vector from the full six-place
box.

The bit-$16$ partner of each base point has center $-x_m$ and the same squared
height and modulus. Conjugation by $J=\operatorname{diag}(1,-1)$ induces
\[
 \begin{pmatrix}a&b\\c&d\end{pmatrix}
 \longmapsto
 \begin{pmatrix}a&-b\\-c&d\end{pmatrix}.
\]
Exact numerator and denominator norm identities prove that this is an
involutive bijection of the two complete bounded candidate sets, including
the nonempty $c=0$ sectors. But $\det J=-1$ is totally negative, so
$J\in\GL_2(\OL)\setminus\GL_2^+(\OL)$ and $J\notin\SL_2(\OL)$; it does not
identify the quotient points. No reflection symmetry of the normalized
quadrilateral is used. All four exact side predicates at all six embeddings
are evaluated separately at each of the twenty centers, and every survivor
set is empty.

The primary rational/SymPy and separate standard-library/Sturm paths agree on
all twenty first-failure histograms and point digests and on all ten labelled
pair aggregates. The primary path additionally freezes the labelled
twenty-mask and ten-pair batch digests
\[
\begin{gathered}
\digest{be161e65208adb27deadd69d0d63cbe5}{a41d500ebd331c02ee7bf22157a47278},\\
\digest{17fcca5c094efb163b7eb39636429b5b}{1b75a56ca6cc4930816742c59607c046}.
\end{gathered}
\]
The primary path replays the already certified v45 root-pair records; the
independent path also reconstructs determinant-one anti-intertwiner bridges
for all twenty labels. Hence the twenty xor-$48$ fibers are pairwise disjoint
and disjoint from the eight proper-order fibers excluded earlier. Removing
these twenty points from the $99$ left by
Proposition~\ref{prop:d7mask1_17} gives $11+68=79$.
\end{proof}


\begin{proposition}[The final two proper-order $d=7$ exclusion pairs]\label{prop:d7finalpairs}
In the canonical point manifest of Theorem~\ref{thm:d7ambient}, none of the
four nontrivial-$R_{83}$ points represented by the free root fibers
\[
 \{13,61\},\quad \{29,45\},\quad \{7,55\},\quad \{23,39\}
\]
lies on the normalized McMullen curve.  The first pair reduces the remainder
of Proposition~\ref{prop:d7batch10} to $77=11+66$; the second reduces it to
\[
 75=11\text{ maximal-multiplier}+64\text{ proper-order}.
\]
\end{proposition}
\begin{proof}
For masks $13$ and $7$, reconstruct the index-$2$ action on the nontrivial
conductor-$83$ ideal lattice, apply the integral determinant-one inversion of
Proposition~\ref{prop:d7mask14}, and then conjugate by
$\operatorname{diag}(q,q^{-1})$.  In the saturated-unit basis
$(\xi_5,\xi_{11},\xi_{13},\xi_{17},4-t^2)$ the exponent vectors are
respectively
\[
 (1,0,1,0,0),\qquad (0,0,0,-1,-2).
\]
Exact multiplication in $\OL$ verifies both units and their integral inverses,
so every reduction and its inverse belongs to $\SL_2(\OL)$.  The reduced
lower-left entries have absolute norm $83^2$.

At each real embedding, rational Sturm isolation supplies intervals
$h_i^-\leq h_i\leq h_i^+$ and $r_i\leq r_i^+$ for the height and modulus.
Before outward decimal rounding, the certificates form the rational upper
endpoints
\[
\begin{aligned}
 \widehat d_i&=\operatorname{upper\_sqrt}(h_i^+/M_i),&
 \widehat c_i&=\operatorname{upper\_sqrt}(1/(M_i h_i^-)),\\
 \widehat n_i&=U_i\widehat d_i,&
 \widehat a_i&=\widehat n_i/h_i^-,\\
 \widehat D_i&=\widehat d_i+\widehat c_i r_i^+,&
 \widehat B_i&=\widehat n_i+\widehat a_i r_i^+.
\end{aligned}
\]
Each stored bound is rounded outward separately from its own hatted endpoint;
no formula substitutes already rounded entries into the last two expressions.
The positive trace form and its exact rational $LDL^{\mathsf T}$ decomposition
therefore enumerate every algebraic integer in each six-place box.  The exact
counts are
\[
\begin{array}{c|rrrrr}
\text{masks}&\#a&\#c&\#(c,d)&c\ne0&c=0\\ \hline
13,29&23917&6625&2043&8768&0\\
7,23 &23889&6609&1955&8514&0
\end{array}
\]
where the last two columns count bounded determinant-one candidates.  For the
second row, the $1955$ bottom rows split into $1954$ rows with $c\ne0$ and the
single row $(c,d)=(0,0)$.  The latter satisfies $ad-bc=0$ for every top row and
therefore contributes no candidate.  This singular sector is part of the
complete enumeration, not discarded by division by $c$.

For each row of the table, bit $16$ changes the center $x$ to $-x$ and leaves
the squared height and modulus fixed.  Conjugation by
$J=\operatorname{diag}(-1,1)$ induces the involution
\[
 \begin{pmatrix}a&b\\c&d\end{pmatrix}
 \longmapsto
 \begin{pmatrix}a&-b\\-c&d\end{pmatrix}
\]
on the complete bounded candidate sets.  Since $\det J=-1$ is totally
negative,
\[
 J\in\GL_2(\OL)\setminus\GL_2^+(\OL),\qquad J\notin\SL_2(\OL).
\]
Thus this involution is only a search bijection and does not identify the
quotient points.  No reflection symmetry of the quadrilateral is assumed:
all four exact circle-side predicates at all six embeddings are evaluated
separately at all four centers.  Every survivor set is empty.

The primary rational/SymPy and separate standard-library/Sturm paths agree on
all counts, failure histograms, and empty survivor sets.  Their point digests
are
\[
\begin{array}{c|l}
13&\digest{d66a04fa445505dd16cf2dad98ee6c9c}{57fe16d90aa4bbeb0b1bf2a1c2cc1a15}\\
29&\digest{b3ce2ca80967f6acf28bf10d6bb1707f}{330c768313680c369dd9d6d2cb3349de}\\
7 &\digest{b653b69bf5888c06df2ee9a20e5f6fd8}{d1d06277de0b862933ba5a4bf8811ae5}\\
23&\digest{392b61bb9ec4dcfd6720de02184032cf}{e211f9094b3e5ed1dac5224be84f8d52}
\end{array}
\]
with pair aggregates
\[
\begin{gathered}
\digest{20ca8ce844583d40514ca6f65db26a30}{2957b057fb0fd23c057cb7e52dab2e68},\\
\digest{aba4002201230fea6eb9e127aa6fa38a}{04e63194d3970adc32c95a1a87386dab}.
\end{gathered}
\]
The xor-$48$ fibers displayed in the statement are mutually disjoint and are
disjoint from all fibers excluded through Proposition~\ref{prop:d7batch10}.
Removing the first two points from $79$ gives $77$, and removing the last two
gives $75$.
\end{proof}

\begin{corollary}[Closure of the nontrivial-$R_{83}$ stratum]\label{cor:d7r83nontrivial}
The exclusions in Propositions~\ref{prop:d7mask14}--\ref{prop:d7finalpairs}
cover the old canonical masks $\{0,1,\ldots,31\}$ exactly once.
Their original-product emptiness and the complete reflected-population
exclusion in Theorem~\externalref{packet}{thm:fullpacketgeometry} exclude all $32$ points
of this stratum. This uses equality of the full oriented populations,
not a correspondence between individual numerical mask labels.
At the historical v53 first-target checkpoint the remaining
$75$ ambient $d=7$ candidates were exactly $11$ maximal-multiplier points,
$32$ proper-$R_{41}$ points, and $32$ principal-Picard-class proper-$R_{83}$
points.
\end{corollary}

\begin{theorem}[Closure of the principal-Picard proper-$R_{83}$ stratum]
\label{thm:d7r83principal}
None of the $32$ ambient quotient points whose multiplier module is in the
principal Picard class of the proper order $R_{83}$ lies on the normalized
McMullen quadrilateral in the trace-polarized
$\mathrm{SL}_2(\OL)$ quotient.
\end{theorem}
\begin{proof}
The v45 ambient classification represents the principal class by the regular
lattice $R_{83}=\OL[\alpha_2]$ with basis $(1,\alpha_2)$. Its $64$
determinant sectors are indexed by masks $0,\ldots,63$. The exact free-root
involution is $m\mapsto m\mathbin{\mathtt{xor}}16$: an integral
determinant-one bridge carries
$\operatorname{Tr}(\alpha_2)I-T_m$ to
$T_{m\mathbin{\mathtt{xor}}16}$. Hence the $32$
disjoint root fibers are represented exactly once by
\[
 \mathcal M=\{0,\ldots,15,32,\ldots,47\}.
\]
The anti-intertwiner used to compare candidate sets has totally negative
determinant. It lies in $\GL_2(\OL)$ but not in $\GL_2^+(\OL)$ or
$\SL_2(\OL)$, so it is not used to identify quotient points. The identity
$\alpha_{23}=-\alpha_2$ gives the same point set for the second conductor-$83$
trace.

For each $m\in\mathcal M$, the certificate constructs and exactly verifies a
cyclotomic unit $q_m$ and applies
$\operatorname{diag}(q_m,q_m^{-1})\in\mathrm{SL}_2(\OL)$.  The six exact fixed-point
heights then give outward rational bounds.  Before rounding, these imply
simultaneous bounds for $c,d$ from the lower-half contribution to the
trace-$LDL^{\mathsf T}$ form and, after $c,d$ are fixed, for $a,b$ from the
upper-half contribution.  Enumerating the resulting finite trace ellipsoids
is therefore complete.

The determinant equation $ad-bc=1$ is solved in both branches.  If $c\ne0$,
the affine solution set is obtained from one exact solution by
$(a,b)=(a_0+nc,b_0+nd)$, and the exact archimedean intervals give all integers
$n$.  If $c=0$, the code separately imposes $ad=1$, integrality of $a=d^{-1}$,
and the complete bounded list of $b$; this essential branch contributes
$49{,}232$ candidates. An image in the original product $P_0$
must satisfy every one of its four closed circle-side inequalities in
all six embeddings, so a failed exact side sign excludes that candidate
from $P_0$. The reflected-product calculation and
Lemma~\externalref{packet}{lem:twochambers}, as assembled in
Theorem~\externalref{packet}{thm:fullpacketgeometry}, supply the full-curve conclusion.

The primary quotient-field/SymPy implementation and a separate
standard-library/Sturm implementation independently reconstruct the order,
root bridges, reductions, bounds, determinant branches, and polygon signs.
For the $32$ centers they agree on $62{,}262$ bounded bottom rows,
$219{,}840$ candidates with $c\ne0$, $49{,}232$ with $c=0$, and hence
$269{,}072$ candidates in total. Every survivor set is empty; all per-center
counts, first-failure histograms, and classification digests agree, with
order-sensitive batch digest
\[
 \digest{7a2f2f6f71efb74627f4f891f108b25}{48a1a9b2a8584cea684eaa630c6d90d90}.
\]
The exact sources are listed in Section~\ref{app:computer}.
\end{proof}

\begin{corollary}[Closure of the proper-$R_{83}$ locus]
\label{cor:d7r83all}
Corollary~\ref{cor:d7r83nontrivial} and
Theorem~\ref{thm:d7r83principal} exhaust and exclude all $64$ proper-$R_{83}$
ambient points. These results and the five earlier maximal decisions
settle $69$ of the $112$ ambient $d=7$ points. Theorem~\ref{thm:maximalallorientations}
below adds the remaining eleven maximal exclusions;
Theorem~\ref{thm:d7r41complete} then excludes all $32$ proper-$R_{41}$
points, completing the $112$-point census.
\end{corollary}

Here ``principal Picard class'' means the identity proper-ideal class of
$R_{83}$; it is unrelated to principal polarization. No new CM-type, Hodge,
or polarization assertion is made in Theorem~\ref{thm:d7r83principal}.



\begin{theorem}[All maximal-multiplier orientations]\label{thm:maximalallorientations}
In the fixed marked extension $M/L$, the maximal-multiplier strata of all
$42$ corrected norm-$\pi$ traces have exactly sixteen ambient quotient
points. Exactly three lie on the normalized McMullen curve: $A,B,C$.
The other thirteen are excluded. This statement concerns maximal
multiplier order, including its occurrences under nonmaximal generated
orders; it makes no proper-order curve-membership assertion.
\end{theorem}
\begin{proof}
Proposition~\ref{prop:alltraceorders} identifies the same maximal-multiplier
population across all $42$ traces; Theorem~\ref{thm:d7ambient} supplies its
sixteen-point census. The exact vertex realizations identify masks
$0,7,15$ with $C,A,B$, respectively. We exclude the thirteen other
masks by checking both targets of Lemma~\externalref{packet}{lem:twochambers}.
First consider the original product $P_0$. Its two historical companion
calculations concern masks $2,13$; the label transfer is justified below.
The eleven remaining original-product searches use masks
\[
 1,3,4,5,6,8,9,10,11,12,14.
\]
Use the maximal action of $\alpha_2=a+bz$, where
$a=-2-t+3t^2-t^4$ and $b=-1+4t+t^2-4t^3+t^5$.
For the established sign representatives $\delta_m$ put
\[
 T_m=\begin{pmatrix}a&-\delta_m b\\ b/\delta_m&a+tb\end{pmatrix}.
\]
The free two-root fibers pair $m$ with $m\mathbin{\mathrm{xor}}16$.
The integral determinant-one matrix
$R_m=\left(\begin{smallmatrix}-1&-\delta_m t\\0&-1\end{smallmatrix}\right)$
satisfies
$R_m(\operatorname{tr}(T_m)I-T_m)=T_{m\mathbin{\mathrm{xor}}16}R_m$.
Thus no negative-determinant reflection is used to identify quotient points.

Balance by $\operatorname{diag}(q,q^{-1})$, with $q=1$ for masks
$1,4,8,9,10$, $q=\xi_{11}^{-1}$ for $3,6$,
$q=\xi_5^{-1}$ for $5,11$, $q=\xi_{17}^{-1}$ for $12$, and
$q=\xi_{13}^{-1}$ for $14$. Here
\[
 \xi_5=t^4-3t^2+1,\quad \xi_{11}=t^3+t^2-2t-1,\quad
 \xi_{13}=t^5-4t^3+3t+1,\quad \xi_{17}=t^3-2t.
\]
All these units and their inverses are integral. With $\eta=\delta_mq^2$,
the balanced fixed tuple has exact invariants
\[
 x=-\eta t/2,\qquad h^2=\eta^2(4-t^2)/4,\qquad |w|^2=\eta^2.
\]
For each real embedding, use the same certified target-polygon bounds
$\operatorname{Im}(\gamma w)\ge M>0$, $|\gamma w|\le U$ as in the
preceding exclusions. For
$\gamma=\left(\begin{smallmatrix}a'&b'\\c'&d'\end{smallmatrix}\right)$
of determinant one,
\[
 |c'w+d'|^2\le h/M,\qquad |a'w+b'|\le U\sqrt{h/M}.
\]
Writing certified rational bounds $h^-\le h\le h^+$ and $|w|\le r^+$,
these imply outward bounds
\[
 |c'|\le\sqrt{1/(Mh^-)},\quad
 |a'|\le U\sqrt{h^+/M}/h^-,\quad
 |d'|\le\sqrt{h^+/M}+|c'|r^+,\quad
 |b'|\le U\sqrt{h^+/M}+|a'|r^+.
\]
Each bound is rounded outward from its own rational pre-rounding endpoint.
The positive trace form gives finite complete containing ellipsoids for
$a'$ and $c'$. For each $c'$, the denominator inequality gives a centered
ellipsoid for $d'$. If $c'\ne0$, the only possible top-right entry is
$b'=(a'd'-1)/c'$; integrality and all numerator bounds are tested exactly.
If $c'=0$, impose $a'd'=1$ and enumerate every $b'$ in the centered
numerator ellipsoid. Thus both determinant branches are exhaustive.

Both implementations independently reconstruct these bounds and enumerate
all candidates, retaining boundary equalities. The final predicates are
necessary polygon-side inequalities in all six embeddings; rational
intervals have exact algebraic sign fallback. Their full results agree:
\begin{center}
\begin{tabular}{rrrr}
\toprule
Mask & Bottom rows & $c'\ne0$ candidates & $c'=0$ candidates\\
\midrule
1 & 1975 & 5280 & 0\\
3 & 1891 & 6996 & 144\\
4 & 2089 & 7736 & 136\\
5 & 1975 & 6748 & 140\\
6 & 2433 & 8540 & 112\\
8 & 1891 & 7392 & 156\\
9 & 2177 & 10840 & 212\\
10 & 2017 & 8956 & 168\\
11 & 1933 & 8972 & 148\\
12 & 2143 & 9796 & 32\\
14 & 2017 & 7004 & 36\\
\midrule
Total & 22541 & 88260 & 1284\\
\bottomrule
\end{tabular}
\end{center}
All $89{,}544$ candidates fail a necessary polygon predicate. The two
portable paths agree on every candidate digest and failure histogram,
with ordered classification digest
\[
 \digest{beff7e2462c7f461f1990c6297381b4}{9c8de2895b59e2be3baa98f0b8d869e2c}.
\]
The full replay covers all eleven masks in both paths. The named exact
certificates are listed in Section~\ref{app:computer}. At this stage
the conclusion is avoidance of $P_0$ for these eleven masks.

For masks $2$ and $13$, the historical companion calculations for
$k=17,38$ and $k=1,22$, respectively, give $9812$ and $9072$
candidates per companion, with none in $P_0$. To transfer these
statements, write $\alpha_k=a_k+b_kz$ and use
\[
 J_k=\begin{pmatrix}1&a_k\\0&b_k\end{pmatrix}.
\]
It intertwines the companion action and multiplication on $(1,z)$.
Here $\OL[\alpha_k]=\OM$, so $b_k$ is an integral unit.
The exact orientation fingerprints select the stated masks and make
$\det(D_{\delta_m}J_k)$ a totally positive unit, where
$D_{\delta_m}=\operatorname{diag}(\delta_m,1)$.
Theorem~\ref{thm:unitnorm} supplies a CM unit whose multiplication
matrix in the companion centralizer has inverse determinant.
Composing with it gives an integral determinant-one intertwiner.
Changing the nonreal generator of this fixed embedded CM action
preserves its unique upper fixed point at every real place.
Thus these are the required labelled $P_0$ statements, without
identifying $\GL_2(\OL)$ with $\SL_2(\OL)$.

For the second target set $R=A_tB_t-B_tA_t$. The two exact target
implementations verify that $\operatorname{tr}(R)=0$, that
$R^2=-\det(R)I$, and that $\det(R)$ is negative at every real
embedding, with absolute norm $21952$. The anti-Mobius action
\[
 r_i(w)=\frac{\sigma_i(R_{11})\overline w+\sigma_i(R_{12})}
                  {\sigma_i(R_{21})\overline w+\sigma_i(R_{22})}
\]
fixes $A_i$ and $B_i$ and is the reflection in their geodesic.
It transforms only the target. The certificates reconstruct the four
reflected vertices and all closed side inequalities independently.
Valid rational lower height and upper modulus bounds for $r_i(Q_i)$,
in the fixed embedding order, are
\[
 10^5(M_i)=(1368,23817,14522,22806,10088,5105),
\]
\[
 10^5(U_i)=(197767,354283,28993,254603,165248,191115).
\]
The exact vertex comparisons imply the bounds on the entire convex
polygon: the horoball $\operatorname{Im}(w)\ge M_i$ and the
geodesic half-plane $|w|\le U_i$ are hyperbolically convex.
The preceding determinant-one bounds therefore give complete finite
ellipsoids for this target too. Certified unit changes of enumeration
coordinates are integral bijections, and every candidate is checked
against the original bounds. Both determinant branches and all boundary
equalities are retained.

For each of the thirteen labels the source module, sign sector and
determinant-one balancing are reconstructed from its frozen exact plan.
The two implementations give the following complete reflected-target
counts, with zero survivors in every row:
\begin{center}
\begin{tabular}{rrrr}
\toprule
Mask & Bottom rows & $c'\ne0$ candidates & $c'=0$ candidates\\
\midrule
1 & 3865 & 83400 & 1332\\
2 & 4277 & 91720 & 1704\\
3 & 3693 & 75752 & 1300\\
4 & 3655 & 81272 & 776\\
5 & 3945 & 81240 & 780\\
6 & 3861 & 89948 & 772\\
8 & 3853 & 88696 & 372\\
9 & 3991 & 97736 & 1800\\
10 & 3819 & 76348 & 280\\
11 & 3945 & 82188 & 708\\
12 & 4143 & 103792 & 1376\\
13 & 4067 & 83148 & 548\\
14 & 4021 & 66496 & 440\\
\midrule
Total & 51135 & 1101736 & 12188\\
\bottomrule
\end{tabular}
\end{center}
In particular all $1{,}113{,}924$ reflected candidates per path fail a
necessary side condition. The canonical preparation, entry lists,
bottom rows, both candidate branches and complete classifications agree
between the separate implementations. The retained executions, their
portable-source correspondence and the exact per-label comparison data
are identified in Section~\ref{app:computer}. These data certify
complete searches; comparison expectations do not define the search
domain. We now have avoidance of both $P_0$ and $P_1$ for each of the
thirteen labelled points. Lemma~\externalref{packet}{lem:twochambers} proves their
full-curve exclusion. Together with the three distinct positive
realizations and the arithmetic census, this proves the theorem.
\end{proof}

\begin{theorem}[Exclusion of the complete proper-$R_{41}$ stratum]
\label{thm:d7r41complete}
In the fixed marked extension $M/L$, none of the $32$ ambient points
with proper multiplier order
\[
 R_{41}=\OL+(41,t-12)\OM
\]
lies on the normalized McMullen curve in the trace-polarized
$\SL_2(\OL)$ quotient.  These are the complete proper-order points
arising from traces $k=14,35$ in Theorem~\ref{thm:d7ambient}.
The distinct embedded order of conductor $(41,t-22)$ is not included
in this exclusion.
\end{theorem}
\begin{proof}
Theorem~\ref{thm:d7nonmax} gives $\operatorname{Pic}(R_{41})=1$
and $N_{M/L}(R_{41}^{\times})=E_L^2$.  Thus the regular module
$R_{41}$ supplies exactly $64$ matrix classes up to $\SL_2(\OL)$ conjugacy,
indexed by the full unit-square masks $0,\ldots,63$.
Write
\[
 \tau=c_{14}=t^5-7t^3-t^2+11t+2,\qquad
 T_m=\begin{pmatrix}0&-\delta_m\pi\\\delta_m^{-1}&\tau\end{pmatrix},
\]
where $\delta_m$ is the binary product in the ordered basis
$(\xi_5,\xi_{11},\xi_{13},\xi_{17},-1,4-t^2)$.
This is the regular basis $(1,\alpha_{14})$, followed by its
orientation twist.  The integral determinant-one matrix
\[
 R_m=\begin{pmatrix}-1&-\delta_m\tau\\0&-1\end{pmatrix}
\]
satisfies
\[
 R_m(\tau I-T_m)=T_{m\mathbin{\mathrm{xor}}16}R_m.
\]
By Lemma~\externalref{packet}{lem:packetpointfibers}, the two root classes are
precisely the fiber over a quotient point.  Hence the canonical masks
\[
 \mathcal M_{41}=\{0,\ldots,15,32,\ldots,47\}
\]
represent all $32$ disjoint pairs $\{m,m\mathbin{\mathrm{xor}}16\}$
exactly once.  The bridge compares the conjugate root action with the
opposite mask; it does not identify actions by a negative-determinant
reflection.  The identity $\alpha_{35}=-\alpha_{14}$ gives the same
point population for the second trace.

For each canonical mask, use the explicit real unit $q_m$ in the
frozen witness table and conjugate by
$\operatorname{diag}(q_m,q_m^{-1})\in\SL_2(\OL)$.
Both exact paths check integrality of the unit and its inverse and
the determinant-one transformation.  With $\eta=\delta_mq_m^2$,
the balanced upper tuple $w=x+ih$, $h>0$, has exact invariants
\[
 x=-\eta\tau/2,\qquad
 h^2=\eta^2(4\pi-\tau^2)/4,\qquad |w|^2=\eta^2\pi.
\]
These identities, total positivity of the squared heights, and the
fixed ordering of the six real embeddings are checked separately
in both implementations.

Here is the finite-domain argument, with the same certified target
bounds as in Theorem~\ref{thm:maximalallorientations}.
At each embedding, a possible image in the normalized quadrilateral
has height at least $M>0$ and modulus at most $U$.
For $\gamma=\left(\begin{smallmatrix}a&b\\c&d\end{smallmatrix}\right)$
of determinant one,
\[
 |cw+d|^2\le h/M,\qquad |aw+b|\le U\sqrt{h/M}.
\]
Certified rational intervals $h^-\le h\le h^+$ and $|w|\le r^+$
therefore give upper bounds
\[
 \widehat D_0\ge\sqrt{h^+/M},\quad
 \widehat C\ge\sqrt{1/(Mh^-)},\quad
 \widehat N_0=U\widehat D_0,\quad
 \widehat A=\widehat N_0/h^-.
\]
The remaining endpoints are
$\widehat D=\widehat D_0+\widehat C r^+$ and
$\widehat B=\widehat N_0+\widehat A r^+$.
Each final entry bound is rounded outward from its own pre-rounding
endpoint.  An integral entry $v$ in the six boxes satisfies
$\Tr_{L/\mathbf Q}(v^2)\le\sum_i R_i^2$.
The positive rational trace form consequently gives finite complete
containing ellipsoids, enumerated by exact $LDL^{\mathsf T}$ recursion.

First enumerate the complete $a$ and $c$ lists.  For each $c$, the
centered denominator ellipsoid around $-cx$, followed by its exact
six-place inequalities, gives all necessary $d$.
For $c\ne0$, each $a$ and bottom row has the unique possible
$b=(ad-1)/c$; test its integrality and all numerator bounds.
For $c=0$, reject $d=0$, impose the unique $a=d^{-1}$, test its
integrality and $a$ bounds, and enumerate every $b$ in the centered
numerator ellipsoid.  Thus the zero branch is included in full.
All boundary equalities remain admissible.

The primary quotient-field path
\path{certificates/d7_r41_v63.py} and the separate standard-library/Sturm
path \path{independent/d7_r41_sturm_v63.py} perform these searches
and the necessary polygon-side tests in all six embeddings.
Their witnesses and exact per-center expectations are retained in
\path{certificates/d7_r41_expected_v63.py}.
The final signs use rational intervals with exact algebraic fallback;
uncertainty is never accepted as an exclusion.  Summed over the
$32$ centers, each implementation obtains
\[
\begin{array}{l|r}
\text{Bounded bottom rows}&62{,}474\\
\text{Candidates with }c\ne0&231{,}744\\
\text{Candidates with }c=0&33{,}468\\
\text{All determinant-one candidates}&265{,}212\\
\text{Survivors}&0.
\end{array}
\]
Every per-center count, candidate digest, zero-branch audit,
first-failure histogram and classification digest agrees.
The ordered classification digest is
\[
 \digest{d3e983df4d9ff3d9d3ebf498df08345b7}{01eb95a5a0edc4a0b5654a88fb4ba06}.
\]
The complete portable replay covers all $32$ masks in each path.
These empty sets prove avoidance of the original product $P_0$.
The full reflected $32$-point search, with fresh target bounds, is part
of Theorem~\externalref{packet}{thm:fullpacketgeometry}. Together with
Lemma~\externalref{packet}{lem:twochambers}, it proves the stated curve exclusion.
\end{proof}

\begin{corollary}[Closure of the nonmaximal $d=7$ geometric packet]
\label{cor:d7geometrycomplete}
The $112$ distinct ambient points attached to the four nonmaximal
$d=7$ traces $k=2,14,23,35$ contain exactly three points of the
normalized McMullen curve: $A,B,C$.  The other $109$ are excluded.
Thus this complete $d=7$ population is classified in the fixed marked,
trace-polarized $\SL_2(\OL)$ model.  Its complement in the full arithmetic census of
Theorem~\externalref{packet}{thm:fullpacketarithmetic} has $256$ points.
\end{corollary}
\begin{proof}
The intrinsic multiplier decomposition of
Theorem~\ref{thm:d7ambient} has $16$ maximal, $64$ proper-$R_{83}$
and $32$ proper-$R_{41}$ points.
Theorem~\ref{thm:maximalallorientations} gives the three maximal
hits $A,B,C$ and thirteen maximal exclusions.
Corollary~\ref{cor:d7r83all} excludes all $64$ proper-$R_{83}$
points, and Theorem~\ref{thm:d7r41complete} excludes the remaining
$32$.  Hence there are $13+64+32=109$ exclusions and three hits.
\end{proof}

\begin{remark}
This closure adds no positive curve point.  The previously proved
geometric CM types and Hodge conclusions at $A,B,C$ retain their
existing scope.  ``Principal Picard class'' for $R_{41}$ refers to
an ideal class and not a principal polarization.  The result concerns
the normalized trace-polarized quotient in the specified marked CM
extension; it neither classifies other CM fields nor identifies the
standard fibers with Mostaed's principally polarized locus.
\end{remark}

\paragraph{Seven pointwise exclusions beyond the $d=7$ population.}
The full arithmetic census also permits individual geometric decisions
outside the complete population of Corollary~\ref{cor:d7geometrycomplete}.
The following result concerns seven specified points, not their whole
multiplier or Picard strata.

\begin{proposition}[Seven further pointwise exclusions]
\label{prop:broadersevenexclusions}
In the fixed marked extension and normalized trace-polarized
$\SL_2(\OL)$ quotient, the seven points in the following table do not
lie on McMullen's curve. The labels are those of the canonical point
manifest of Theorem~\externalref{packet}{thm:fullpacketarithmetic}; all have orientation
mask zero. The last three columns count the complete bounded matrices
before the original-product closed polygon tests, in either exact
implementation. These counts are retained historical $P_0$ data.
\[
\begin{array}{l|r|l|r|r|r}
\text{Point label}&k&\text{Multiplier conductor}&c\ne0&c=0&\text{Total}\\\hline
\texttt{R127:pic0:00}&10&(127,t-94)&6572&2304&8876\\
\texttt{R43\_20:pic0:00}&4&(43,t-20)&7704&952&8656\\
\texttt{R7sq:pic0:00}&12&(7,t-1)^2&7084&1320&8404\\
\texttt{R41\_22:pic0:00}&5&(41,t-22)&7148&1616&8764\\
\texttt{R43\_31:pic0:00}&8&(43,t-31)&8556&1436&9992\\
\texttt{R7:pic0:00}&0&(7,t-1)&8840&588&9428\\
\texttt{R127:pic1:00}&10&(127,t-94)&8488&0&8488.
\end{array}
\]
The survivor count is zero in every row. Here \texttt{pic0} means the
principal Picard module class, not a principal polarization;
\texttt{pic1} is the nonprincipal class represented at conductor
$(127,t-94)$ by $J_9$ in the proof of
Theorem~\externalref{packet}{thm:fullpacketarithmetic}.
\end{proposition}

\begin{proof}
Write $\alpha_k=A_k+B_kz$ as in the arithmetic census. For the six
principal rows use the ordered basis $(1,B_kz)$, and for the last row
use $(1+9z,B_{10}z)$. Put $s=0$ or $9$ accordingly and
$n=1+st+s^2$. Multiplication by $\alpha_k$ has the integral matrix
\[
 T_{k,s}=\begin{pmatrix}
 A_k-sB_k&-B_k^2\\
 n&A_k+(t+s)B_k
 \end{pmatrix}.
\]
The two paths independently reconstruct this action, its characteristic
polynomial and its upper fixed tuple $w=x+ih$. In the fixed ordering of
the six real embeddings its exact invariants are
\[
 x=-\frac{B_k(t+2s)}{2n},\qquad
 h^2=\frac{B_k^2(4-t^2)}{4n^2},\qquad
 |w|^2=\frac{B_k^2}{n}.
\]
All squared heights are positive. The conductor, module and root-fiber
identifications are those of Theorem~\externalref{packet}{thm:fullpacketarithmetic}.
In particular the last point has root partner
$(J_{100},16)$; this is the other root representative of the same point,
not an additional exclusion.

The portable witness data specify a balancing transformation
$G\in\SL_2(\OL)$, consisting of a diagonal unit transformation,
with inversion first when prescribed. Both paths verify its integral
inverse and reconstruct the transported point. If $p'=Gp$, an integral
determinant-one representative $H$ at $p'$ corresponds to $HG$ at $p$.
Thus this change of representative preserves the exact orbit question.

Apply the finite-domain argument of
Theorem~\ref{thm:d7r41complete} to the transported tuple. It uses only
positive source heights, determinant one and the certified target
height/modulus bounds, and therefore applies to the nonprincipal row
as well. Its outward rational bounds and positive trace form contain
every possible entry in finite ellipsoids. The exact unit changes of
enumeration coordinates are integral bijections: every enumerated entry
is mapped back before the original six-place bounds are tested.

For clarity, write an original shifted bound at embedding $i$ as
$|v_i+q_i x_i|^2+q_i^2h_i^2\le L_i^2$.
For an integral unit $u$, choose certified rational
$B_i\ge |u_i|L_i$ and set $y=uv$.
Then every admissible $v$ satisfies
\[
 \Tr_{L/\mathbf Q}((y+uqx)^2)
 \le \sum_i B_i^2-\Tr_{L/\mathbf Q}((uq)^2h^2).
\]
This is the searched ellipsoid centered at $-uqx$;
$q=0$ gives the entry-box searches, $q=c$ the denominator search,
and $q=a$ the zero-branch numerator search.
Multiplication by $u^{-1}$ recovers every original integral entry,
regardless of the embedding signs of $u$.
Strictly negative remaining radii give empty domains; zero radii and
boundary equality are retained.

For every admitted $c$, the denominator enumeration finds all admitted
$d$. If $c\ne0$, each admitted $a$ has the unique possible completion
$b=(ad-1)/c$; integrality and the original numerator inequalities are
tested. If $c=0$, require $d\ne0$ and $a=d^{-1}\in\OL$, check the
$a$ bounds and enumerate the complete remaining $b$ domain. Hence the
zero branch is included even when it returns no matrix.

For the last row the portable computation also rejects nonprimitive
bottom rows: $(c,d)\ne\OL$ precludes $ad-bc=1$. The ideal index is
computed afresh; a unit $c$ directly certifies primitivity. On every
remaining nonzero bottom row it verifies integer identities
\[
 M_cR=D M_d,\qquad M_c v=-D e_0,\qquad D>0,
\]
where $M_c,M_d$ are multiplication matrices in the integral power basis
and $e_0$ represents $1$. These identities certify the complete affine
formula $b=(Ra+v)/D$ before scanning the entire original $a$ list.
Simultaneous negation preserves determinant, action and every original
bound. Both matrix signs are restored before canonical sorting and
polygon classification; the retained matrices receive exact determinant
checks. These reductions therefore omit no bounded candidate.

The six principal rows give $54{,}120$ complete candidates per path.
The last row freshly gives $29{,}555$ possible $a$ entries, $8{,}167$
possible $c$ entries, $2{,}005$ bottom rows and $8{,}488$ complete
candidates. The primary exact algebra/trace-ellipsoid path and the
separate integer/Sturm path agree on the reconstructed preparation,
entry lists, bottom rows, both candidate branches and complete polygon
classifications. The named portable certificates and comparison data
are listed in Section~\ref{app:computer}. Comparison targets do not
define or truncate the search; the full success status requires
recomputed counts and all eight canonical artifact hashes to agree.

Every candidate fails a necessary closed polygon-side condition in at
least one embedding, with exact algebraic signs or certified rational
intervals. This proves avoidance of $P_0$ for these seven labels.
Their separate reflected searches are included in the full $365$-label
certificate of Theorem~\externalref{packet}{thm:fullpacketgeometry};
Lemma~\externalref{packet}{lem:twochambers} then proves the seven curve exclusions.
\end{proof}

\begin{corollary}[Historical subpacket coverage]\label{cor:geometry119}
The complete $d=7$ population and the seven additional labels above form
a $119$-point subset of the $368$-point arithmetic census: the three
vertices and $116$ exclusions. Its complement has $249$ points,
all excluded by Theorem~\externalref{packet}{thm:fullpacketgeometry}. The figures
$119$ and $249$ describe the historical division of the first-target
work, not a current coverage total or an unresolved set.
\end{corollary}
\begin{proof}
The $112$-point $d=7$ census and the seven distinct additional points
are disjoint. Outside this subset the four proper $d=3$ strata
contribute $128-4=124$ points, conductor $(41,t-22)$ contributes
$32-1=31$, and conductor $(127,t-94)$ contributes $96-2=94$.
Their sum is $249$. Principal controls are not counted a second time,
and a paired root is not a second quotient point. The first-target
coverages $116$ and $249$ are disjoint and exhaust all $365$ nonvertices.
\end{proof}

\begin{remark}
These seven negative orbit decisions add no positive geometric CM type,
$K$-signature, Hermitian discriminant or Hodge-algebra assertion. They
neither exclude a whole remaining conductor or Picard stratum nor
identify any standard fiber with a principally polarized model or with
Mostaed's named principally polarized locus. The geometric CM and Hodge
conclusions at $A,B,C$ are unchanged.
\end{remark}

